% Einheit 3 von 5 – Dreieck (bank/flaechen/e3.jsonl, zusammenbau v0.1)
%% TODO Zweigzeile Teil 1: was hier gelernt wird (Satz in Schülersprache aus dem Einheitstitel) – Einheit 3
\einheitenkopf[e3]{Einheit 3 von 5 · Dreieck}
\zweigzeile{ab Kl. 7, je nach Buch bis Kl. 8 (am Gymnasium ab Kl. 5) · P10 oft}
\setcounter{aufgabe}{17}

%% TODO Titel: Ich-kann-Satz fehlt in der Bank (Platzhalter: Kettenname) – Nr. 18
%% TODO Anweisung: ein Satz über der ersten Teilaufgabe fehlt in der Bank – Nr. 18
\begin{aufgabe}{Dreieck}
%% TODO \rechenplatz{n}: Schrittzahl des Grundfalls fehlt in der Bank (nur bei fester Schreibform, 2.3 b)
\begin{teile}
\teil Zeichne die Höhe zur Grundseite $g$ ein und markiere den rechten Winkel.

\dreieck{(0,0)}{(5,0)}{(2,3)}{}{}{$g$}{}{}{}
\teil Dreieck, $g = 8$ cm, $h = 7$ cm – Fläche? A = \leerfeld[cm²]
\teil Rechtwinkliges Dreieck – Fläche?

\dreieckrw{5}{2.67}{$15$ cm}{$8$ cm}{$17$ cm}
A = \leerfeld[cm²]
\teil Stumpfwinkliges Dreieck, Grundseite $6$ cm; die Höhe dazu liegt außerhalb und ist $4$ cm lang – Fläche?

\dreieck{(0,0)}{(3,0)}{(4.5,2)}{}{}{$6$ cm}{}{}{}
A = \leerfeld[cm²]
\teil Dreieck, $g = 4{,}6$ cm, $h = 3$ cm – Fläche? A = \leerfeld[cm²]
\teil Dreieck, $A = 26$ cm², $h = 4$ cm – Grundseite? g = \leerfeld[cm]
\teil Ein gleichseitiges Dreieck hat die Seitenlänge $s$. Term für den Umfang? u = \leerfeld
\teil Ein dreieckiges Grundstück ABC hat bei C einen rechten Winkel; $BC = 16$ m, $AC = 45$ m. Ein Teilstück QBC mit Q auf AC soll $120$ m² groß sein. Wie lang ist QC? \hfill (P10 2020 OS) \\ QC = \leerfeld[m]
\end{teile}
\end{aufgabe}

%% TODO Titel: Ich-kann-Satz fehlt in der Bank (Platzhalter: Kettenname) – Nr. 19
%% TODO Anweisung: ein Satz über der ersten Teilaufgabe fehlt in der Bank – Nr. 19
\begin{aufgabe}{Höhe zur passenden Grundseite wählen (drei Höhen)}
\begin{teile}
\teil Dreieck ABC: $AB = 8$ cm, $BC = 7{,}5$ cm. Die Höhe auf AB ist $6$ cm, die Höhe auf BC $6{,}4$ cm. Fläche – rechne mit BC und der passenden Höhe. A = \leerfeld[cm²]
\end{teile}
\end{aufgabe}

%% TODO Titel: Ich-kann-Satz fehlt in der Bank (Platzhalter: Kettenname) – Nr. 20
%% TODO Anweisung: ein Satz über der ersten Teilaufgabe fehlt in der Bank – Nr. 20
\begin{aufgabe}{Umfang}
\begin{teile}
\teil Dreieck mit den Seiten $5$ cm, $7$ cm und $9$ cm – Umfang? u = \leerfeld[cm]
\end{teile}
\end{aufgabe}

%% TODO Titel: Ich-kann-Satz fehlt in der Bank (Platzhalter: Kettenname) – Nr. 21
%% TODO Anweisung: ein Satz über der ersten Teilaufgabe fehlt in der Bank – Nr. 21
\begin{aufgabe}{Dreieck – Fehler finden, Begründen, Anwendung, Darstellungswechsel}
\begin{teile}
\teil Dreieck mit $g = 9$ cm und $h = 6$ cm. Jana rechnet: \rechnung{A &= 9 \cdot 6 \\ A &= 54 \text{ cm}^2} Wo ist der Fehler?
\teil Warum rechnet man beim Dreieck Grundseite mal Höhe und nimmt davon die Hälfte? Begründe mit einem Parallelogramm.
\teil Ein dreieckiger Hausgiebel ist unten $8{,}4$ m breit und $3{,}5$ m hoch. Er wird gestrichen; eine Dose Farbe reicht für $6$ m². Wie viele Dosen braucht man?
\teil Zeichne ein Dreieck, dessen Umfang $u = 3 \cdot k$ ist, und beschrifte die Seiten.

\rechenplatz[halb]{4}
\end{teile}
\end{aufgabe}
